%% parts/p02/propositional-metatheory.tex --- P2.1 metatheory (Phase 4)
%% Statements only.  Proofs are in annex/p02/propositional-metatheory.tex.
\chapter{Propositional Metatheory}
\label{sec:p02-propositional-metatheory}

This chapter keeps the semantic, natural-deduction, sequent, and resolution
standings distinct.  The exact inherited natural-deduction rules in
\cref{sec:prop-deduction-01} and classical-sequent rules in
\cref{sec:prop-deduction-02} govern the corresponding statements below.  The
registry sheets \hyperref[sheet:nd-cpl]{\texttt{nd-cpl}} and
\hyperref[sheet:lk]{\texttt{lk}} are added navigation and comparison
apparatus; they do not replace those inherited rule displays.  The added
resolution development uses
\hyperref[sheet:resolution-cpl]{\texttt{resolution-cpl}}.  The broader
metatheory treatment is \cite{EbbinghausFlumThomas1994,OpenLogic2026}; the
repaired resolution attribution is recorded below.

\section{Soundness}
\label{sec:p02-prop-soundness}

\begin{hbtheorem}{thm:p02-nd-soundness}{proved=annex:p2.3.1}{Soundness of classical propositional natural deduction}
For every \(\Gamma\subseteq\mathsf{Fm}_{\mathrm{CPL}}\) and every
\(\varphi\in\mathsf{Fm}_{\mathrm{CPL}}\),
\[
 \Gamma\derives_{\mathsf{ND}_{\mathrm{CPL}}}\varphi
 \quad\text{implies}\quad
 \Gamma\entails\varphi.
\]
\end{hbtheorem}

\begin{hbtheorem}{thm:p02-lk-soundness}{proved=annex:p2.3.2}{Soundness of the classical propositional sequent calculus}
Every sequent derivable in the propositional restriction of
\(\mathsf{LK}\) is satisfied by every valuation, with satisfaction understood
exactly as in the sequent definition in \cref{sec:prop-deduction-02}.
\end{hbtheorem}

The conclusion says that a derivable sequent is \emph{satisfied} by every
valuation.  It does not call a sequent ``valid'' as though it were a formula;
that distinction preserves repair R-022.

\section{Maximal consistent sets and countervaluations}
\label{sec:p02-prop-maximal-consistency}

\begin{hbdefinition}{def:p02-consistency}{Consistency, propositional}
A set \(\Gamma\subseteq\mathsf{Fm}_{\mathrm{CPL}}\) is
\emph{consistent} iff
\(\Gamma\notderives_{\mathsf{ND}_{\mathrm{CPL}}}\bottom\).
\end{hbdefinition}

\begin{hbdefinition}{def:p02-maximal-consistent}{Maximal consistent set, propositional}
A \emph{maximal consistent set} is a consistent
\(\Delta\subseteq\mathsf{Fm}_{\mathrm{CPL}}\) such that
\(\Delta\cup\{\varphi\}\) is inconsistent for every
\(\varphi\notin\Delta\).
\end{hbdefinition}

\begin{hblemma}{thm:p02-mcs-closure}{proved=annex:p2.3.3}{Closure properties of maximal consistent sets}
If \(\Delta\) is maximal consistent, then, for all formulas
\(\varphi,\psi\):
\begin{enumerate}
  \item if \(\Delta\derives_{\mathsf{ND}_{\mathrm{CPL}}}\varphi\), then
        \(\varphi\in\Delta\);
  \item exactly one of \(\varphi\) and \(\neg\varphi\) belongs to
        \(\Delta\);
  \item \(\varphi\land\psi\in\Delta\) iff
        \(\varphi\in\Delta\) and \(\psi\in\Delta\);
  \item \(\varphi\lor\psi\in\Delta\) iff
        \(\varphi\in\Delta\) or \(\psi\in\Delta\); and
  \item \(\varphi\to\psi\in\Delta\) iff
        \(\varphi\notin\Delta\) or \(\psi\in\Delta\).
\end{enumerate}
\end{hblemma}

\begin{hblemma}{thm:p02-lindenbaum}{proved=annex:p2.3.4}{Countable propositional Lindenbaum lemma}
Every consistent subset of the fixed countable language
\(\mathsf{Fm}_{\mathrm{CPL}}\) is contained in a maximal consistent set.
\end{hblemma}

\begin{hblemma}{thm:p02-truth-lemma}{proved=annex:p2.3.5}{Propositional truth lemma}
Let \(\Delta\) be maximal consistent and define
\(v_{\Delta}(p_n)=1\) iff \(p_n\in\Delta\).  Then, for every
\(\varphi\in\mathsf{Fm}_{\mathrm{CPL}}\),
\[
 v_{\Delta}\entails\varphi
 \quad\text{iff}\quad
 \varphi\in\Delta.
\]
\end{hblemma}

\begin{hbmethod}{def:p02-maximal-countervaluation}{Maximal-consistent-set countervaluation}
To refute an underivable consequence, first prove
\(\Gamma\cup\{\neg\varphi\}\) consistent, extend it along the fixed
enumeration of formulas, and read a valuation from membership of the
propositional variables.  Lemma~\ref{thm:p02-truth-lemma} then makes every
member of \(\Gamma\) true and \(\varphi\) false.  This method depends on the
nonvacuous-discharge convention in \cref{sec:prop-deduction-01};
the annex records the required no-occurrence case separately.
The name abbreviates this three-stage proof method, not a new semantic
relation; the completeness proof below uses precisely these stages.
\end{hbmethod}

\section{Completeness and compactness}
\label{sec:p02-prop-completeness}

\begin{hbtheorem}{thm:p02-nd-completeness}{proved=annex:p2.3.6}{Strong completeness of classical propositional natural deduction}
For every \(\Gamma\subseteq\mathsf{Fm}_{\mathrm{CPL}}\) and every
\(\varphi\in\mathsf{Fm}_{\mathrm{CPL}}\),
\[
 \Gamma\entails\varphi
 \quad\text{implies}\quad
 \Gamma\derives_{\mathsf{ND}_{\mathrm{CPL}}}\varphi.
\]
\end{hbtheorem}

Together with Theorem~\ref{thm:p02-nd-soundness}, this gives
\(\Gamma\derives_{\mathsf{ND}_{\mathrm{CPL}}}\varphi\) iff
\(\Gamma\entails\varphi\) for the fixed countable language.

\begin{hbtheorem}{thm:p02-compactness}{proved=annex:p2.3.7}{Propositional compactness, fixed countable language}
For every \(\Gamma\subseteq\mathsf{Fm}_{\mathrm{CPL}}\) and every formula
\(\varphi\):
\begin{enumerate}
  \item if \(\Gamma\entails\varphi\), some finite
        \(\Gamma_0\subseteq\Gamma\) satisfies
        \(\Gamma_0\entails\varphi\); and
  \item \(\Gamma\) is satisfiable iff every finite subset of \(\Gamma\) is
        satisfiable.
\end{enumerate}
\end{hbtheorem}

\begin{hbtheorem}{thm:p02-arbitrary-prop-compactness-bpi}{cited=HowardRubin1998}{Boolean Prime Ideal Theorem and unrestricted propositional compactness}
Working over \(\mathsf{ZF}\), the Boolean Prime Ideal Theorem (equivalently,
the Ultrafilter Lemma of Definition~\ref{def:ultrafilter-lemma}) is equivalent
to the following scheme: for every set \(P\) of propositional variables and
every set \(\Gamma\) of finitary formulas over \(P\), if each finite subset
of \(\Gamma\) is satisfiable, then \(\Gamma\) is satisfiable.
\end{hbtheorem}

See also \cite{Jech1973Choice,PaseauLeek2022} for the choice-principle and
logical formulations of this equivalence.

Theorem~\ref{thm:p02-compactness} is weaker in set-theoretic strength: its
fixed countable alphabet supports the explicit enumeration used in
Lemma~\ref{thm:p02-lindenbaum}, with no appeal to a Boolean prime ideal.
The equivalence above is the treatment required by D-P1-005; it does not
silently promote the source proof to arbitrary languages.

\hbstub{D-P2-011: add the full arbitrary-language syntax and completeness
development, including exact choice dependencies beyond the compactness
scheme stated above.}

\begin{hbtheorem}{thm:p02-lk-completeness}{sketched=annex:p2.3.8}{Completeness of the propositional sequent calculus}
For all finite multisets \(\Gamma,\Delta\) of propositional formulas, the
sequent \(\Gamma\Rightarrow\Delta\) is derivable in the propositional
restriction of \(\mathsf{LK}\) iff every valuation making every member of
\(\operatorname{supp}(\Gamma)\) true makes some member of
\(\operatorname{supp}(\Delta)\) true.
\end{hbtheorem}

The scope includes empty antecedents, empty succedents, and multiple-formula
succedents.  The retained source sketch uses Cut and does not claim cut-free
completeness; cut elimination belongs to Part~7.

\section{Interpolation}
\label{sec:p02-prop-interpolation}

\begin{hbdefinition}{def:p02-interpolant}{Craig interpolant, propositional}
If \(\varphi\entails\psi\), a \emph{Craig interpolant} for that consequence
is a formula \(\theta\) such that
\(\varphi\entails\theta\), \(\theta\entails\psi\), and every variable of
\(\theta\) occurs in both \(\varphi\) and \(\psi\).  The name is needed to
distinguish this formula-level witness from an arbitrary intermediate
consequence: its common-vocabulary restriction is the nonredundant clause used
by the theorem and construction below.
\end{hbdefinition}

\begin{hbtheorem}{thm:p02-interpolation}{cited=Craig1957}{Propositional interpolation theorem}
Every semantic consequence \(\varphi\entails\psi\) in classical
propositional logic has a Craig interpolant.
\end{hbtheorem}

\begin{hbmethod}{def:p02-interpolant-truth-table}{Truth-table construction of a propositional interpolant}
Let \(X\) be the finite set of variables common to \(\varphi\) and \(\psi\),
ordered by variable index.
For each assignment \(u:X\to\{0,1\}\) extendible to a valuation satisfying
\(\varphi\), form the complete conjunction of \(X\)-literals describing
\(u\), and disjoin these conjunctions in lexicographic order with \(0<1\).  Empty
conjunction and disjunction use \(\topv\) and \(\bottom\), respectively.
The resulting formula uses only variables from \(X\); its two consequence
clauses are stated separately below so that this derived construction is not
attributed to Craig's original proof.
Part~7 gives the proof-theoretic Maehara construction.
\end{hbmethod}

\begin{hblemma}{thm:p02-interpolant-truth-table-correctness}{proved=annex:p2.3.9}{Correctness of the truth-table interpolant construction}
If \(\varphi\entails\psi\), the output of
Method~\ref{def:p02-interpolant-truth-table} is a Craig interpolant for that
consequence.
\end{hblemma}

\section{Resolution}
\label{sec:p02-prop-resolution}

\textbf{OPEN:} This section owes a reviewed citable destination for the exact
source definition of propositional literals and clauses at source lines
1949--1982.  That definition is present earlier in Part~II but intentionally
unlabeled; this section does not invent a Definition alias.  Clause formation
and inference remain separate operations.  The only inference rule is registered on
\hyperref[sheet:resolution-cpl]{\texttt{resolution-cpl}}.

\begin{hbdefinition}{def:p02-resolvent}{Resolvent, propositional}
Let \(E,F\) be clauses, let \(L\in E\), and let the complementary literal
\(\overline L\) belong to \(F\).  The \emph{resolvent of \(E,F\) on \(L\)}
is
\[
 (E\setminus\{L\})\cup(F\setminus\{\overline L\}).
\]
A \emph{resolution derivation} from a clause set \(S\) is a finite sequence
whose members belong to \(S\) or are resolvents of earlier members.
The term names a derived clause, not a formula conversion or a new connective;
that distinction is required by the soundness and completeness statements
below.
\index{resolution derivation}
\end{hbdefinition}

\begin{hbtheorem}{thm:p02-resolution-soundness}{proved=annex:p2.3.10}{Soundness of propositional resolution}
Every valuation satisfying both parent clauses satisfies their resolvent.
Consequently, if the empty clause has a resolution derivation from \(S\),
then \(S\) is not satisfiable.
\end{hbtheorem}

\begin{hbtheorem}{thm:p02-resolution-completeness}{proved=annex:p2.3.11}{Refutational completeness of propositional resolution}
For every set \(S\) of propositional clauses, \(S\) is not satisfiable iff
the empty clause has a resolution derivation from \(S\).
\end{hbtheorem}

\begin{hbmethod}{def:p02-davis-putnam}{Davis--Putnam variable elimination}
For a finite clause set \(S\), first delete its tautological clauses.  If a
variable occurs in what remains, choose one such variable \(p\), retain the
clauses mentioning neither \(p\) nor \(\neg p\), and add each non-tautological
resolvent between a \(p\)-parent and a \(\neg p\)-parent.  Call the resulting
set \(S_p\).  The variable \(p\) does not occur in \(S_p\), so \(S_p\) has
strictly fewer
occurring variables than the tautology-free remainder; and \(S_p\) is
satisfiable iff \(S\) is.  This
elimination step is due to Davis and Putnam \cite{DavisPutnam1960}; Robinson's
later resolution calculus is cited separately \cite{Robinson1965}.  Finiteness
is consumed by the variable-count induction, while
Theorem~\ref{thm:p02-compactness} extends refutational completeness to
arbitrary clause sets.
\end{hbmethod}

\begin{figure}[tp]
\centering
\begin{tikzpicture}[
  clause/.style={draw,rounded corners,inner sep=4pt,font=\small},
  every edge/.style={draw,-{Latex[length=2mm]}},
  level distance=13mm,
  sibling distance=31mm]
\node[clause] (c1) at (0,0) {\(\{P,Q\}\)};
\node[clause] (c2) at (3,0) {\(\{\neg P,R\}\)};
\node[clause] (c4) at (6,0) {\(\{\neg R\}\)};
\node[clause] (c3) at (9,0) {\(\{\neg Q,R\}\)};
\node[clause] (np) at (4.5,-1.5) {\(\{\neg P\}\)};
\node[clause] (nq) at (7.5,-1.5) {\(\{\neg Q\}\)};
\node[clause] (q) at (2.25,-3) {\(\{Q\}\)};
\node[clause] (empty) at (4.875,-4.5) {\(\varnothing\)};
\draw[-{Latex[length=2mm]}] (c2) -- node[left,font=\scriptsize] {\(R\)} (np);
\draw[-{Latex[length=2mm]}] (c4) -- (np);
\draw[-{Latex[length=2mm]}] (c3) -- node[right,font=\scriptsize] {\(R\)} (nq);
\draw[-{Latex[length=2mm]}] (c4) -- (nq);
\draw[-{Latex[length=2mm]}] (c1) -- node[left,font=\scriptsize] {\(P\)} (q);
\draw[-{Latex[length=2mm]}] (np) -- (q);
\draw[-{Latex[length=2mm]}] (q) -- node[left,font=\scriptsize] {\(Q\)} (empty);
\draw[-{Latex[length=2mm]}] (nq) -- (empty);
\end{tikzpicture}
\caption{Resolution-refutation DAG.  Each two-parent merge is an instance of
the rule on \hyperref[sheet:resolution-cpl]{\texttt{resolution-cpl}}; the last node is the empty clause
(apparatus A-121).}
\label{fig:p02-resolution-dag}
\end{figure}

\section{Decision procedures and complexity}
\label{sec:p02-prop-decision-complexity}

\begin{hbtheorem}{thm:p02-decidability}{cited=EbbinghausFlumThomas1994}{Decidability of classical propositional logic}
There is a terminating algorithm which, on finite encoded inputs, decides
whether a propositional formula is satisfiable or valid, whether two formulas
are logically equivalent, and whether a finite set \(\Gamma\) semantically
entails a formula \(\varphi\).
\end{hbtheorem}

\begin{hbconvention}{def:p02-complexity-standing}{Complexity standing for propositional decision problems}
A decision problem is here a set \(L\subseteq\{0,1\}^*\).  The class
\(\mathsf P\) consists of the decision problems decided by a deterministic
multitape Turing machine in a number of steps bounded by a polynomial in the
input length.  A polynomial-time
many-one reduction from \(L\) to \(K\) is a polynomial-time computable map
\(f\) such that \(x\in L\) iff \(f(x)\in K\).  The class \(\mathsf{NP}\)
consists of those \(L\) for which some polynomial \(q\) and polynomial-time
predicate \(R\) satisfy, whenever \(n\) is the encoded length of \(x\),
\[
 x\in L\quad\text{iff}\quad
 \text{some string }y\text{ of length at most }q(n)\text{ satisfies }R(x,y).
\]
The class \(\mathsf{coNP}\) consists of complements of \(\mathsf{NP}\)
problems.  For a complexity class \(\mathcal C\), a problem \(K\) is
\(\mathcal C\)-complete under these reductions iff \(K\in\mathcal C\) and
every \(L\in\mathcal C\) has a polynomial-time many-one reduction to \(K\).
These metalinguistic definitions are stated locally because the complexity
claims use the distinctions now; Part~4 supplies the full class ladder.
\index{P@\(\mathsf P\)}\index{NP@\(\mathsf{NP}\)}
\index{coNP@\(\mathsf{coNP}\)}\index{polynomial-time many-one reduction}
\end{hbconvention}

\begin{hbdefinition}{def:p02-sat-taut}{SAT and TAUT}
Write each variable \(p_n\) as the character \texttt{p}, the ordinary binary
digits for \(n\), and the terminator \texttt{;}; use the fully parenthesized formula string;
and encode each character from the resulting finite alphabet by a fixed-width
binary block.  This metalinguistic encoding and its inverse are linear-time
and linearly intertranslatable with ordinary finite-alphabet formula encodings.
\(\mathsf{SAT}\) is the set of encodings of satisfiable formulas, and
\(\mathsf{TAUT}\) is the set of encodings of valid formulas.
Malformed bit strings belong to neither language; syntactic recognition is
linear-time for the fixed encoding.
These metalinguistic abbreviations are needed to distinguish existential
satisfiability from universal validity.  They are not the same decision
problem, although negation relates their complements, and both distinctions
are used in the next two results.
\index{SAT}\index{TAUT}
\end{hbdefinition}

\begin{hbtheorem}{thm:p02-sat-np-complete}{cited=AroraBarak2009}{SAT NP-completeness}
\(\mathsf{SAT}\) is \(\mathsf{NP}\)-complete under polynomial-time many-one
reductions.
\end{hbtheorem}

Cook's theorem-proving formulation is the historical source
\cite{Cook1971}; the status citation supports the modern formulation above.

\begin{hbcorollary}{thm:p02-taut-conp-complete}{cited=AroraBarak2009}{TAUT is coNP-complete}
\(\mathsf{TAUT}\) is \(\mathsf{coNP}\)-complete under polynomial-time
many-one reductions.
\end{hbcorollary}

The preceding classification does not settle whether
\(\mathsf P=\mathsf{NP}\) or \(\mathsf{NP}=\mathsf{coNP}\).  Part~4 locates
these open equalities in the complexity-class ladder; Part~12 treats SAT
solvers as algorithms rather than as complexity-class collapses.

\begin{hbdefinition}{def:p02-tractable-classes}{Horn, dual-Horn, 2-CNF, and XOR instances}
A clause is \emph{Horn} when it has at most one positive literal and
\emph{dual-Horn} when it has at most one negative literal.  A clause set is
\emph{2-CNF} when every clause has at most two literals.  An
\emph{XOR instance} is a finite system of linear equations over the
two-element field, presented as parity constraints on Boolean variables.
For a CNF input with no empty clause, let \(N\) denote its total number of
literal occurrences; an empty clause is detected first and makes the instance
unsatisfiable.  For a parity input, let \(m\) and \(n\) be its numbers of
equations and variables.
The corresponding satisfiability problems restrict their inputs to these
classes.  The four names distinguish nonidentical syntactic restrictions whose
different closure patterns license the algorithms stated below; the table uses
the distinction rather than treating them as synonyms for ``easy SAT.''
\index{Horn clause}\index{dual-Horn clause}\index{2-CNF}\index{XOR-SAT}
\end{hbdefinition}

\begin{hbtheorem}{thm:p02-horn-sat}{cited=DowlingGallier1984}{Horn-SAT tractability}
Horn satisfiability is decidable in time linear in the input size by forward
chaining (equivalently, unit propagation in the implicational presentation):
after the empty-clause check, the bound is \(O(N)\).
\end{hbtheorem}

\begin{hbcorollary}{thm:p02-dual-horn-sat}{proved=annex:p2.3.12}{Dual-Horn-SAT tractability}
Dual-Horn satisfiability is decidable in \(O(N)\) time.
\end{hbcorollary}

\begin{hbtheorem}{thm:p02-two-sat}{cited=AspvallPlassTarjan1979}{2-SAT tractability}
For a 2-CNF instance without an empty clause, form the implication graph on
literals by adding directed edges from \(\overline L\) to \(K\) and from
\(\overline K\) to \(L\) for every clause \(L\lor K\), treating a unit clause
\(L\) as \(L\lor L\).  The instance
is satisfiable iff no occurring variable \(p\) has both \(p\) and \(\neg p\)
in one strongly connected component.  The test runs in \(O(N)\) time.
\end{hbtheorem}

\begin{hbtheorem}{thm:p02-xor-sat}{cited=Schaefer1978}{XOR-SAT tractability}
XOR satisfiability is decidable by Gaussian elimination over the two-element
field, for example in \(O(mn\min\{m,n\})\) field operations with a dense
implementation.
\end{hbtheorem}

\begin{table}[tp]
\centering
\small
\begin{tabular}{@{}p{0.18\textwidth}p{0.37\textwidth}p{0.28\textwidth}@{}}
\toprule
Class & Recognition and decision method & Worst-case status \\
\midrule
Horn & At most one positive literal per clause; forward chaining & \(O(N)\) \\
Dual-Horn & At most one negative literal; complement variables and use Horn & \(O(N)\) \\
2-CNF & Clauses of size at most two; no literal shares an SCC with its complement & \(O(N)\) \\
XOR & Linear parity equations over the two-element field; Gaussian elimination & \(O(mn\min\{m,n\})\) dense field operations \\
\bottomrule
\end{tabular}
\caption{Tractable propositional satisfiability classes (apparatus A-030).}
\label{tab:p02-tractable-sat}
\end{table}

\section{Failure modes and further reading}
\label{sec:p02-prop-failure-modes}

\begin{hbwarning}{def:p02-propositional-ambiguity}{Propositional ambiguity}
Without a separate justification, do not reverse an implication, read
disjunction exclusively, infer a disjunct from a disjunction, affirm the
consequent, omit a variable from a truth table, cite soundness in the
completeness direction, or infer satisfiability merely because a resolution
search has not yet produced the empty clause.
\end{hbwarning}

For detailed classical semantics and proof systems, see
\cite{EbbinghausFlumThomas1994,OpenLogic2026,Kleene1952}; for resolution and
normal-form transformations, see
\cite{Robinson1965,BaazEglyLeitsch2001}; and for complexity, see
\cite{Cook1971,AroraBarak2009}.  These references form the propositional
contribution to apparatus A-186; the Part-2 finalizer will combine it with the
first- and higher-order reading lists.
